\documentclass[11pt,a4paper]{article}
% =====================================================================
%  Shared preamble for the Part V lecture notes (lectures 19-22).
%
%  These four lectures sit outside Geron, so the notes ARE the primary
%  source and are examined on the same terms as the textbook parts.
%  Everything here is chosen to make that readable on paper: one column,
%  generous measure, and the same six-colour palette as the slides so a
%  student moving between the two is not re-learning what red means.
%
%  Build with pdflatex (twice, for the toc and the refs):
%      cd notes && latexmk -pdf lecture-19.tex
%  or  make -C notes
% =====================================================================
\usepackage[T1]{fontenc}
\usepackage[utf8]{inputenc}
\usepackage{newtxtext,newtxmath}      % Times-like: prints well, reads well
\usepackage[scaled=0.86]{inconsolata} % code, at a size that sits beside Times
\usepackage{microtype}
% newtx defines \Bbbk, \openbox, \proof and \endproof; amssymb and amsthm
% then redefine them and abort. Freeing the four names before those two
% packages load is the standard fix, and it costs nothing here: we use
% amsthm only for \newtheorem.
\let\Bbbk\relax
\usepackage{amsmath,amssymb}
\let\openbox\relax
\let\proof\relax
\let\endproof\relax
\usepackage{amsthm}
\usepackage{booktabs}
\usepackage{enumitem}
\usepackage{xcolor}
\usepackage{tcolorbox}
\usepackage{titlesec}
\usepackage{graphicx}
\usepackage[hidelinks]{hyperref}
\tcbuselibrary{skins,breakable}

% ---- the course palette, from assets/css/custom.css -------------------
\definecolor{aimlprimary}{HTML}{0B3D62}   % deep blue  - structure
\definecolor{aimlaccent} {HTML}{C0392B}   % brick red  - failures
\definecolor{aimlsuccess}{HTML}{14663A}   % green      - repairs
\definecolor{aimlmath}   {HTML}{6C3483}   % purple     - the derivation
\definecolor{aimlmuted}  {HTML}{4B5563}
\definecolor{aimlrule}   {HTML}{B0BCC7}
\definecolor{aimlsoft}   {HTML}{F4F7F9}

\usepackage[a4paper,margin=32mm,top=30mm,bottom=30mm]{geometry}
\setlength{\parskip}{0.45\baselineskip}
\setlength{\parindent}{0pt}
\linespread{1.06}

% ---- headings --------------------------------------------------------
\titleformat{\section}{\Large\bfseries\color{aimlprimary}}{\thesection}{0.7em}{}
\titleformat{\subsection}{\large\bfseries\color{aimlprimary}}{\thesubsection}{0.7em}{}
\titleformat{\subsubsection}{\bfseries\color{aimlmuted}}{}{0em}{}
\titlespacing*{\section}{0pt}{1.6\baselineskip}{0.5\baselineskip}
\titlespacing*{\subsection}{0pt}{1.2\baselineskip}{0.35\baselineskip}

% ---- boxes -----------------------------------------------------------
% One box per role, and the role is the colour. A student who has seen the
% slides already knows what each one means before reading a word of it.
\newtcolorbox{keybox}[1][]{
  breakable, enhanced, colback=aimlsoft, colframe=aimlprimary,
  boxrule=0.4pt, left=8pt, right=8pt, top=6pt, bottom=6pt,
  fonttitle=\bfseries\small, coltitle=white, colbacktitle=aimlprimary,
  attach boxed title to top left={yshift=-2mm, xshift=4mm},
  boxed title style={boxrule=0pt, arc=1pt}, #1}

\newtcolorbox{failbox}[1][]{
  breakable, enhanced, colback=white, colframe=aimlaccent,
  boxrule=0.9pt, left=8pt, right=8pt, top=6pt, bottom=6pt,
  fonttitle=\bfseries\small, coltitle=white, colbacktitle=aimlaccent,
  attach boxed title to top left={yshift=-2mm, xshift=4mm},
  boxed title style={boxrule=0pt, arc=1pt}, #1}

\newtcolorbox{derivbox}[1][]{
  breakable, enhanced, colback=white, colframe=aimlmath,
  boxrule=0.6pt, left=8pt, right=8pt, top=6pt, bottom=6pt,
  fonttitle=\bfseries\small, coltitle=white, colbacktitle=aimlmath,
  attach boxed title to top left={yshift=-2mm, xshift=4mm},
  boxed title style={boxrule=0pt, arc=1pt}, #1}

% An exercise. Deliberately the same shape as the notebook's `try` field:
% one change, and what should happen to the output.
\newtcolorbox{trybox}[1][]{
  breakable, enhanced, colback=white, colframe=aimlsuccess,
  boxrule=0.5pt, left=8pt, right=8pt, top=6pt, bottom=6pt,
  fonttitle=\bfseries\small, coltitle=white, colbacktitle=aimlsuccess,
  attach boxed title to top left={yshift=-2mm, xshift=4mm},
  boxed title style={boxrule=0pt, arc=1pt}, #1}

% ---- small semantics -------------------------------------------------
\newcommand{\scope}[1]{\textcolor{aimlmuted}{\small #1}}
\newcommand{\term}[1]{\textbf{#1}}
\newcommand{\code}[1]{\texttt{\small #1}}
\newcommand{\lecture}[1]{Lecture~#1}

\newtheorem{definition}{Definition}[section]
\newtheorem{proposition}[definition]{Proposition}

% ---- title block -----------------------------------------------------
% \notestitle{19}{Information retrieval: the lexical foundation}{SciFact (BEIR)}
\newcommand{\notestitle}[3]{%
  \thispagestyle{empty}
  \begin{flushleft}
    {\color{aimlmuted}\small\sffamily
     APPLICATIONS OF MACHINE LEARNING \quad$\cdot$\quad
     BSc Mathematics of Artificial Intelligence}\\[2mm]
    {\color{aimlrule}\rule{\textwidth}{0.8pt}}\\[4mm]
    {\color{aimlmuted}\large Lecture #1 \quad$\cdot$\quad Part V \quad$\cdot$\quad
     Extended notes}\\[3mm]
    {\Huge\bfseries\color{aimlprimary}\baselineskip=1.15\baselineskip #2\par}\vspace{4mm}
    {\color{aimlmuted}\large Dataset: #3}\\[3mm]
    {\color{aimlrule}\rule{\textwidth}{0.8pt}}
  \end{flushleft}
  \vspace{2mm}
  \begin{keybox}[title={Why these notes exist}]
    Lectures 19--22 sit outside G\'eron, the course's primary text. For those
    four lectures \emph{these notes are the primary source}, and they are
    examined on exactly the same terms as the textbook parts. Everything in
    the slide deck for this lecture is here, with the arguments written out at
    length rather than compressed onto a slide; the numbers are the ones the
    accompanying notebook prints, and every one of them is a measurement on
    the corpus named above rather than a figure quoted from a paper.
  \end{keybox}
  \vspace{1mm}
  {\small\tableofcontents}
  \vspace{2mm}
  \noindent{\color{aimlrule}\rule{\textwidth}{0.4pt}}
  \vspace{1mm}
}


\begin{document}
\notestitlefor{1}{I}{What machine learning is, and how we will work}{California housing}{Chapters 1--2}

\section{What the course is about}

The course has a one-sentence description: how to build a machine learning
system, and how to know whether it works. The second clause is the hard one.
Fitting a model is a few lines of library code, and it has been for a decade.
Establishing that the number the model reports means what it appears to mean is
the whole of the rest of the course, and it is the part that does not get
easier with better libraries.

\subsection{Why the second clause is hard}

A bug in a training loop does not crash. It returns a plausible number. That
sentence is the reason this course exists in the shape it does, and it is worth
five concrete instances before we go any further:

\begin{itemize}[leftmargin=1.4em,itemsep=2pt]
  \item a scaler fitted before the train/test split --- information leaks from
        the test set into the training statistics, the score improves, and the
        system is broken in deployment;
  \item a missing \code{model.eval()} --- dropout stays active while you
        evaluate, so the reported accuracy is of a model nobody will ever ship;
  \item a missing \code{optimizer.zero\_grad()} --- gradients accumulate across
        steps, and nothing anywhere raises an error;
  \item a metric averaged per batch rather than over the set --- wrong by
        construction whenever the batches are not the same size;
  \item accuracy on an imbalanced set --- 90\% from a model that has learned to
        predict one class.
\end{itemize}

\begin{failbox}[title={The property that organises this course}]
Every one of those five runs to completion. Every one of them prints a number.
Not one raises an exception, and four of the five produce a number that is
\emph{better} than the honest one, which is precisely the direction in which
you are least likely to go looking for a bug.
\end{failbox}

Compilers catch type errors and test suites catch crashes. Neither catches a
number that is merely wrong, and machine learning produces those in quantity.
What catches them is evidence: a baseline beside the number, an ablation
showing the number moves when you remove the component that supposedly earned
it, and a stated failure case. Those three are the working habits this course
is trying to install, and they are examined directly.

\subsection{What kind of problem needs learning at all}

A learning system is worth building when the rule you would otherwise write by
hand is long, unstable, or unknown. Spam filtering is the standard example:
the hand-written version is a list of patterns that its adversary edits around
every week, so the maintenance never ends. A learning system replaces the list
with a procedure for deriving the list, and the maintenance becomes retraining.

Three consequences follow, and they are worth stating before any method:

\begin{itemize}[leftmargin=1.4em,itemsep=2pt]
  \item \term{If a short exact rule exists, write it.} A learning system is
        strictly worse than \code{if amount > limit} when that is genuinely the
        rule: it is slower, larger, harder to audit, and occasionally wrong.
  \item \term{The system is only as good as what it learns from.} There is no
        method in this course that repairs data which does not contain the
        answer, and a great deal of the course is about noticing that case.
  \item \term{It will fail quietly.} See above. This is why the evaluation
        material is not an appendix.
\end{itemize}

\subsection{The taxonomy, briefly}

The distinctions are standard, and the reason to have them straight is that
they determine which methods are even candidates.

\begin{center}
\small
\begin{tabular}{lp{0.62\textwidth}}
\toprule
\textbf{Supervision} & \textbf{What the training data carries} \\
\midrule
supervised      & every instance carries its expected output \\
unsupervised    & no labels; structure has to come from the inputs alone \\
semi-supervised & a few labelled instances among many unlabelled \\
self-supervised & labels manufactured from the data itself --- mask a word,
                  predict it \\
reinforcement   & no labels; a reward signal for a sequence of actions \\
\bottomrule
\end{tabular}
\end{center}

Cutting a different way: \term{instance-based} learning keeps the training
examples and answers a new query by similarity to them; \term{model-based}
learning fits parameters and discards the examples. And cutting a third way:
\term{batch} learning trains once on everything held in memory, \term{online}
learning consumes a stream and updates as it goes.

\subsection{The failure modes, named up front}

Four, and every one of them will reappear with a measurement attached:

\begin{itemize}[leftmargin=1.4em,itemsep=2pt]
  \item \term{Insufficient data.} The method is not the bottleneck.
  \item \term{Unrepresentative data.} The training set is not a sample of the
        population the system will meet. This is the one that is invisible in
        every metric you compute on your own data, and \S\ref{sec:split} is
        about buying insurance against it.
  \item \term{Overfitting.} The model has learned the training set, including
        its noise. Low training error, high error on anything else.
  \item \term{Underfitting.} The model is too rigid to represent the pattern
        that is there. Both errors high.
\end{itemize}

The only way to distinguish the last two --- indeed the only way to know how a
system generalises at all --- is to try it on cases it has never seen. Hence a
training set and a test set, hence the rule that the test set is spent once,
and hence the fact that the split happens before the exploration rather than
after it.

\section{The working method}

You will write machine learning code with an assistant: if not in this course,
then in the thesis after it and the job after that. This course is explicit
about how, rather than pretending otherwise, and it spends the first half of
Lecture 1 on it because the alternative is spending the rest of the course
undoing the habits.

The argument for spending the time is short. Applied machine learning is a
small amount of library glue and an enormous amount of data work and
evaluation. The scarce skill is judgement about whether a result can be
trusted. An assistant supplies the glue and none of the judgement.

\subsection{The four rules}

\begin{keybox}[title={Four rules for AI-assisted work}]
\begin{enumerate}[leftmargin=1.6em,itemsep=2pt]
  \item \term{Never keep code you cannot explain line by line.}
  \item \term{Every number needs a baseline.} A metric with nothing to compare
        it to is decoration.
  \item \term{Every model needs an ablation.} Remove a component; show the
        number moves.
  \item \term{Report the failure cases.} A system with no known failure mode
        has not been tested.
\end{enumerate}
\end{keybox}

Rules 2 to 4 are what a result has to come with before anyone should believe
it --- yours, a colleague's, or an assistant's. They are also, roughly, Part C
of the written paper. Rule 1 is Part B: the paper puts a cell in front of you
and asks what would have to be true for its number to be trusted.

\subsection{The loop, and whose job each step is}

\begin{center}
\small
\begin{tabular}{clp{0.50\textwidth}c}
\toprule
\textbf{\#} & \textbf{Step} & \textbf{What it means} & \textbf{Who} \\
\midrule
1 & Specify & state the input, the output, the constraint, the check & you \\
2 & Generate & and then stop & the assistant \\
3 & Read & as a reviewer, not as the author & you \\
4 & Test & against a case whose answer you know & you \\
5 & Verify & that the number means what you think it means & you \\
\bottomrule
\end{tabular}
\end{center}

Four of the five are yours. The one that is not is the one that was never the
hard part. Steps 3 to 5 are, quite literally, the syllabus.

\subsubsection{Step 1 --- Specify}

A usable specification names four things: the \term{input} (shape, types,
provenance), the \term{output} (exactly what object, with what guarantees), the
\term{constraint} (what must not happen), and the \term{check} (how you will
know it worked). Written out for a task from later in this lecture:

\begin{keybox}[title={A specification, in full}]
\textbf{Task:} prepare the housing features for a model.\\[2mm]
\term{input} --- the 16{,}512-row training frame, 8 numeric columns and 1
categorical, with 168 missing values in \code{total\_bedrooms}\\
\term{output} --- a fitted \code{ColumnTransformer} and the transformed
matrix\\
\term{constraint} --- the imputer and the scaler are fitted on the training
rows \emph{only}; the test rows are transformed with those statistics and never
contribute to them\\
\term{check} --- the output has 16{,}512 rows and 13 columns (8 numeric plus 5
one-hot levels) and no \code{NaN}
\end{keybox}

The part people omit is the third. Almost every silent failure in machine
learning is a missing constraint, not a missing feature --- and the constraint
above is exactly the one whose violation produces the leaked scaler from the
list at the top of these notes.

Every code cell in every notebook in this course is preceded by a box in that
form. The box is step 1; the cell below it is what step 2 returned; the check
line is step 4, written down in advance. Read the box, answer the check in your
head, then run the cell, and you have done the whole loop with only the free
step handed to you.

\subsubsection{Step 2 --- Generate, then stop}

The dangerous moment is not the generation. It is the reflex to run it.
Generated code is read \emph{before} it is executed, because output you have
already seen is very hard to disbelieve. In a domain where wrong code still
returns a plausible number, running first destroys the only cheap opportunity
to notice.

\subsubsection{Step 3 --- Read it like a reviewer}

Five questions, in this order, every time:

\begin{enumerate}[leftmargin=1.6em,itemsep=2pt]
  \item \term{What touched the test set?} Trace every line backwards to its
        data source.
  \item \term{What was fitted, and on what?} \code{fit} and \code{transform}
        are different verbs.
  \item \term{What is the shape here?} Silent broadcasting is silent.
  \item \term{What was dropped?} Rows, columns, \code{NaN}s --- count them.
  \item \term{What is the default I did not ask for?} Every unnamed argument
        is a decision someone else made.
\end{enumerate}

Question 5 is not pedantry. Later in this lecture a single unrequested default
--- a colormap --- produces a figure that is colourful, ran without error, and
is wrong.

\subsubsection{Step 4 --- Test against a case you already know}

Not ``does it run''. Does it give the right answer on an input whose answer you
can state independently? Four kinds of case, all cheap:

\begin{itemize}[leftmargin=1.4em,itemsep=2pt]
  \item \term{A shape.} 16{,}512 rows in, 16{,}512 rows out --- and if one-hot
        encoding added five columns, 13 out, not 9.
  \item \term{A count.} 168 missing values before, 0 after.
  \item \term{A degenerate case.} A constant column has zero variance. What
        does the scaler do with it?
  \item \term{Arithmetic you can do on paper.} A layer with 32 inputs and 32
        outputs has $32 \times 32 + 32 = 1{,}056$ parameters. Print the number
        and compare.
\end{itemize}

\begin{failbox}[title={The rule that separates testing from watching}]
If you cannot state the expected answer \emph{before} running the cell, you are
not testing. You are watching.
\end{failbox}

\subsubsection{Step 5 --- Verify that the number means what it appears to}

The code is correct and the number is still misleading. This is the step nobody
does, and it is four questions:

\begin{itemize}[leftmargin=1.4em,itemsep=2pt]
  \item \term{Against what?} What does the trivial model score?
  \item \term{On which rows?} Training, validation or test --- and were they
        chosen before or after anything was fitted?
  \item \term{How much does it move?} A difference smaller than the spread
        across splits has not been measured.
  \item \term{In what units?} An error of 70{,}000 is meaningless until you
        know the prices run from 15{,}000 to 500{,}000.
\end{itemize}

Every one of the four is a question about evidence, not about code.

\section{The brief}

You are a data scientist at a real-estate investment company. The task: use
California census data to build a model of housing prices. The data gives
population, median income and median housing price for each block group --- the
smallest geographical unit the US Census publishes, which we will call a
\term{district}. The model predicts a district's median housing price from the
other metrics.

\subsection{The first question is never technical}

What is the business objective? Building a model is not a goal. How the company
uses the output determines the framing, the algorithm, the metric, and how much
effort is worth spending.

The answer here: the prediction is fed to \emph{another} machine learning
system, together with other signals, to decide whether a district is worth
investing in. A piece of information fed to a downstream system is called a
\term{signal}, after Shannon.

\begin{keybox}[title={The consequence of being a component}]
Nobody downstream ever looks at your prediction. A component that quietly
degrades inside a chain of components is not noticed by anyone --- which is
exactly why the evidence has to come from you, unprompted, rather than from a
complaint.
\end{keybox}

\subsection{What the current solution is}

District prices are currently estimated manually by experts: a team gathers
information and applies complex rules. When they can check their estimates
against the truth, they are typically off by more than 30\%.

Read that figure carefully, because two things about it matter later. It is a
\emph{relative} criterion --- 30\% of the district's price --- and in the next
lecture we choose a metric measured in dollars, which is not the same thing. It
is also the stakeholder's claim about their own team, not a measurement we
made. We will treat it as an assumption and label it as one wherever it
appears.

\subsection{Framing, in four questions}

\begin{center}
\small
\begin{tabular}{llp{0.44\textwidth}}
\toprule
\textbf{Question} & \textbf{Answer} & \textbf{Because} \\
\midrule
Supervision   & supervised & every district carries its expected output \\
Task          & multiple, univariate regression & predict one value, from
                several features \\
Learning mode & batch, offline & no data stream, and 20{,}640 rows fit in
                memory \\
\bottomrule
\end{tabular}
\end{center}

``Multiple'' refers to the several input features; ``univariate'' to the single
output. Predict several values per district and it becomes multivariate
regression; add a data stream and it becomes online learning.

The fourth question is the one that can sink the project: \emph{what is
everyone assuming?}

\begin{failbox}[title={The assumption that would have sunk this project}]
Suppose the downstream system converts our prices into categories --- cheap,
medium, expensive --- and discards the numbers. Then this is a
\emph{classification} problem, every regression metric we are about to choose
is the wrong one, and we would have built the wrong system competently.

We asked. They need the actual prices. We may proceed --- and the value of the
question is not that the answer was reassuring, but that it was cheap and the
alternative was months.
\end{failbox}

\section{Looking at the data, once}

Ten attributes, one row per district. The habit is \code{head()} then
\code{info()}, always, before anything else.

\subsection{What \code{info()} reports}

\code{RangeIndex: 20640 entries}, ten columns, nine of them \code{float64} and
one \code{object}. Two facts fall straight out of it.

\subsubsection{207 districts are incomplete}

\code{total\_bedrooms} has 20{,}433 non-null values against 20{,}640 rows.
Three options, all legitimate: drop those districts, drop the whole attribute,
or fill the gaps with a chosen value. We return to this in the build, and the
choice is not obvious --- dropping 207 rows is cheap, dropping a column that
correlates with the target is not, and filling requires a decision about
\emph{what} to fill with and \emph{on which rows the filling value is
computed}. That last clause is the constraint from \S2.

\subsubsection{One attribute is not numerical}

\code{ocean\_proximity} takes repeated values from a small set: this is a
categorical attribute, not free text.

\begin{center}
\small
\begin{tabular}{lr}
\toprule
\textbf{Category} & \textbf{Districts} \\
\midrule
\code{<1H OCEAN}  & 9{,}136 \\
\code{INLAND}     & 6{,}551 \\
\code{NEAR OCEAN} & 2{,}658 \\
\code{NEAR BAY}   & 2{,}290 \\
\code{ISLAND}     & 5 \\
\bottomrule
\end{tabular}
\end{center}

Note \code{ISLAND}: five districts in the whole state. Rare categories cause
trouble later --- a one-hot column that is nonzero five times in 20{,}640 rows
can appear in a training fold and not in a validation fold, and a stratified
split cannot protect a stratum that small.

\subsection{Four things the histograms show}

Plotting all nine numeric attributes at fifty bins takes one line and repays it
immediately.

\subsubsection{Observation 1 --- the income is not in dollars}

\code{median\_income} runs from roughly 0.5 to 15. That is not a salary. The
data was scaled and capped: at 15 (in fact 15.0001) at the top and 0.5 (0.4999)
at the bottom. The units are roughly tens of thousands of dollars, so 3 means
about \$30{,}000. Preprocessed attributes are normal and not a defect;
understanding how they were computed is not optional.

\subsubsection{Observation 2 --- the target is capped}

\code{median\_house\_value} is capped, and it is our \emph{label}.

\begin{failbox}[title={Why a capped label is more serious than a capped feature}]
The model will learn that prices never exceed the cap, because in the training
data they never do. If the client needs accurate predictions above
\$500{,}000, the model is not merely inaccurate there --- it is structurally
incapable of producing such a number.

Two honest fixes: collect proper labels for the capped districts, or remove
them from \emph{both} the training and the test set. Removing them from the
training set alone would leave you scoring a model on districts it was
forbidden to learn about.
\end{failbox}

The spike is not small: 992 districts, nearly 5\% of the data, carry a label
that is a ceiling rather than a price.

\subsubsection{Observation 3 --- the scales differ wildly}

\code{total\_rooms} runs from 2 to 39{,}320; \code{median\_income} from 0.4999
to 15.0001. Whether that matters depends entirely on the method.

\begin{center}
\small
\begin{tabular}{p{0.44\textwidth}p{0.44\textwidth}}
\toprule
\textbf{Scaling is required for} & \textbf{Scaling does nothing for} \\
\midrule
gradient descent, regularised linear models (ridge, lasso), SVMs,
$k$-nearest neighbours, $k$-means, PCA, neural networks
&
decision trees, random forests, ordinary least squares \\
\bottomrule
\end{tabular}
\end{center}

The left column is anything that measures a distance, penalises a coefficient,
or takes a step. The right column has reasons: a tree splits on one feature at
a time, so any monotone rescaling produces the identical tree; and least
squares is equivariant under an invertible affine map of the features --- the
fit is identical and the coefficients simply absorb the change.

All three models built in Lecture 2 are in the right-hand column, and that is
not an incidental remark. It is why the scale-before-split leak, which is the
canonical beginner's error, will cost \emph{nothing} measurable here. Lecture 2
measures it over twenty seeds rather than asserting it.

\subsubsection{Observation 4 --- many attributes are right-skewed}

They extend much further to the right of the median than to the left, which
makes patterns harder for some algorithms to detect. The usual repair is a
transform toward symmetry: a square root or a logarithm. A feature following a
power law is a good candidate for $\log$ --- districts of 10{,}000 people are
about ten times rarer than districts of 1{,}000, not exponentially rarer, and
the log of such a variable is roughly linear in its rank.

\section{The split, before the exploration}\label{sec:split}

Four observations from four histograms, and we have not yet plotted one
variable against another. That is exactly the moment to stop.

\subsection{Data snooping is a bias, not a discipline problem}

Your brain is an outstanding pattern detector, which means it overfits. If you
study the test set, you will --- without any dishonesty, and without being able
to point at the moment it happened --- choose a model that suits it. Your
estimate of generalisation error will be optimistic and you will ship something
that underperforms.

The defence is procedural rather than moral: create the test set \emph{before}
exploring any further, and do not look at it again. Typically 80/20, though the
proportion is a function of size --- with ten million instances, holding out
1\% is plenty.

\subsection{Why random sampling is not automatically enough}

Purely random sampling is fine if the dataset is large relative to the number
of attributes. Ours may not be, and random sampling then risks \term{sampling
bias}: a test set whose composition differs from the population, so that the
error measured on it is an error on the wrong population.

Surveyors do not call 1{,}000 people at random; they ensure the sample is
representative on the dimensions that matter. The experts told us median income
predicts price, so the test set must be representative across income levels.
That is \term{stratified sampling}: bin the stratifying variable, then sample
within each bin in proportion.

The bins have to be chosen so that no stratum is too small to sample from. Five
categories, cut at 1.5, 3.0, 4.5 and 6.0, leave the smallest --- category 1 ---
with 822 districts, which is comfortably enough.

\subsection{Does it actually matter? Measure it}

The point of this course is that the previous paragraph is a hypothesis until
someone counts. Comparing each split's income-category proportions against the
full dataset's:

\begin{center}
\small
\begin{tabular}{lr}
\toprule
\textbf{Split} & \textbf{Largest error in stratum proportion} \\
\midrule
purely random & up to 6.4\% \\
stratified    & up to 0.36\% \\
\bottomrule
\end{tabular}
\end{center}

Roughly eighteen times better, for one keyword argument. Note the shape of the
argument: the stratification was justified by an expert's claim, and then the
claim was checked with a measurement rather than deferred to.

\subsection{Assert after every structural step}

\begin{keybox}[title={Three lines, three real bugs}]
\begin{verbatim}
assert len(strat_train_set) + len(strat_test_set) == 20640
assert set(strat_train_set.columns) == set(strat_test_set.columns)
assert strat_train_set.index.intersection(strat_test_set.index).empty
\end{verbatim}
Nothing was dropped; the two frames have the same columns; no district is in
both. Each of the three has been a real bug in a real notebook.
\end{keybox}

An assertion is a claim about your data that fails \emph{loudly}, and in a
domain that otherwise fails silently that makes it the most valuable line in
the file. Write one after every split, merge, reshape and drop, and count what
you dropped. Reviewer question 4 is then answered by the program instead of by
you, every time it runs, including the times you forget to ask.

\section{Exploring --- the training set, and only that}

From here everything is computed on the 16{,}512-district training split. The
other 4{,}128 are sealed until the last slide of Lecture 2. Work on a copy, so
that nothing invented while exploring --- a new column, a filter, a dropped row
--- can reach the split frame by accident, and so that every number below is
honestly labelled as a training-set number.

\subsection{Look at it, and tune the plot until you can}

The data is geographical, so plot latitude against longitude. It looks like
California, and beyond that nothing: at full opacity the ink saturates and
every district is the same black dot.

Setting \code{alpha=0.2} changes the picture completely --- the Bay Area, Los
Angeles, San Diego and the Central Valley appear, because now the density is
encoded in the darkness. Same data, one keyword argument. Expect to tune
plotting parameters until the pattern shows; a plot that shows nothing is
usually a plot, not an absence.

Adding two more channels --- radius for population, colour for price --- shows
that prices relate strongly to location and to population density. Red on the
coast, blue inland. The model will not know what ``coast'' means, but it will
learn longitude.

\subsection{A default nobody chose, and how to catch it}

The three-channel plot above is conventionally drawn with \code{cmap="jet"}. It
ran, it produced something colourful, and it is wrong.

A colour scale is a claim: equal steps in the data look like equal steps on the
screen. That claim is testable --- plot each colormap's lightness against
position along it.

\begin{center}
\small
\begin{tabular}{lrrl}
\toprule
\textbf{Colormap} & \textbf{Lightness range} & \textbf{Largest reversal} &
\textbf{Verdict} \\
\midrule
\code{jet}     & 83.0 & 1.06 & rises, falls, rises \\
\code{turbo}   & 78.9 & 0.87 & better, still not monotone \\
\code{viridis} & 75.9 & 0.00 & every step is a step \\
\bottomrule
\end{tabular}
\end{center}

A non-monotone lightness profile invents bands the data does not have and
buries detail in the yellow. For the roughly 8\% of men with a red--green
deficiency, \code{jet} reverses rank order outright. Photocopy it and it stops
being a scale at all.

\begin{failbox}[title={Notice the shape of this failure}]
The code ran. Nothing errored. The defect was in a default nobody chose
deliberately, and it took a measurement to see. That is every failure in this
course, and it is only the first lecture.
\end{failbox}

\subsection{Correlations, and what they cannot see}

Pearson's $r$ against the target, on the training split:

\begin{center}
\small
\begin{tabular}{lr}
\toprule
\textbf{Attribute} & \textbf{$r$ with \code{median\_house\_value}} \\
\midrule
\code{median\_income}      & 0.688380 \\
\code{total\_rooms}        & 0.137455 \\
\code{housing\_median\_age}& 0.102175 \\
\code{households}          & 0.071426 \\
\code{total\_bedrooms}     & 0.054635 \\
\code{population}          & $-0.020153$ \\
\code{longitude}           & $-0.050859$ \\
\code{latitude}            & $-0.139584$ \\
\bottomrule
\end{tabular}
\end{center}

Income dominates: one bar is five times the next. That is the feature we
already stratified on, on the experts' word, before we had looked --- and the
measurement says they were right.

These are training-set values. Compute them on all 20{,}640 rows and the third
decimal moves, which is why a number is worth nothing without the set it was
measured on attached to it.

\begin{failbox}[title={Two things $r$ does not tell you}]
$r$ measures \emph{linear} correlation only. A dataset can have $r = 0$ and be
completely dependent --- any symmetric nonlinear relationship will do it. And
$r = \pm 1$ says nothing whatever about the slope: your height in inches
correlates $1.0$ with your height in nanometres.
\end{failbox}

\subsection{The cap, made visible --- and the lines that are not there}

Plot the strongest predictor against the target and the \$500{,}001 cap appears
as a hard horizontal line: 787 districts in our training split. There are
fainter horizontal lines too, and this is where the lecture's method earns its
keep. A horizontal stripe means many districts share one exact price. That is a
property of the table, not of California, so it is \emph{countable} --- and
squinting at a scatter plot is not counting.

A well-known description of this dataset names three fainter lines. Counted in
our own training split:

\begin{center}
\small
\begin{tabular}{lrl}
\toprule
\textbf{Value} & \textbf{Districts} & \textbf{Verdict} \\
\midrule
\$350{,}000 & 62 & real --- 21$\times$ the background \\
\$450{,}000 & 31 & marginal \\
\$280{,}000 & 3  & indistinguishable from background \\
\bottomrule
\end{tabular}
\end{center}

A typical price is shared by 3 districts, so the third is not a line at all.
Meanwhile the five commonest values below the cap --- \$137{,}500,
\$162{,}500, \$112{,}500, \$187{,}500 and \$225{,}000 --- go unmentioned in
that description, and they are the strongest stripes in the plot. Every one
sits on a multiple of \$12{,}500, which tells you something about how the
survey recorded prices and nothing about the housing market.

\begin{keybox}[title={The habit, not the numbers}]
Three values were asserted; one survives. Counting took one line. You will
inherit descriptions of data for the rest of your career --- from papers, from
colleagues, from an assistant --- and the cost of checking one against the data
is almost always smaller than the cost of building on it.
\end{keybox}

\subsection{Attribute combinations}

Raw counts per district are nearly meaningless because districts differ in
size. Ratios are not.

\begin{center}
\small
\begin{tabular}{lr}
\toprule
\textbf{Attribute} & \textbf{$r$ with price} \\
\midrule
\code{median\_income}   & 0.688 \\
\code{rooms\_per\_house}& 0.144 \\
\code{total\_rooms}     & 0.137 \\
\code{bedrooms\_ratio}  & $-0.256$ \\
\bottomrule
\end{tabular}
\end{center}

\code{bedrooms\_ratio} --- bedrooms as a fraction of all rooms --- is far more
informative than either raw count it is built from, and its sign is
interpretable: a district whose rooms are mostly bedrooms is a district of
small dwellings.

\begin{failbox}[title={A warning you will need next lecture}]
When you create combined features, check that they are not too linearly
correlated with the existing ones, and in particular avoid simple weighted sums
of features already present. \term{Collinearity} causes real problems for some
models --- linear regression above all. Lecture 2 derives the normal equation
and shows exactly why: the matrix you have to invert becomes singular, or close
enough to it that the inverse is numerically meaningless.
\end{failbox}

\section{The notebook}

The Lecture 1 notebook runs on Colab's free CPU tier in about two minutes and
contains everything above in order: the fetch, \code{info()}, the missing
values, the histograms, the stratified split, and the exploration. Every code
cell carries the specification that would produce it.

\begin{trybox}[title={What to change}]
Run it top to bottom once. Then change one thing --- the number of income bands
in \code{pd.cut} --- and watch the sampling-bias table. Two lines, and it is
the whole argument of \S\ref{sec:split}: too few strata and the stratification
buys you little; too many and each stratum is too small for its own proportion
to be estimated reliably, so the cure starts to resemble the disease.
\end{trybox}

\section{Exercises}

Five, in the style of the written paper. Solutions on Lecture 2's deck.

\begin{enumerate}[leftmargin=1.6em,itemsep=4pt]
  \item A colleague reports that their model reaches an RMSE of \$48{,}000 on
        the California housing data, and that they chose the model by trying
        six of them and keeping the best. State, in one sentence, what that
        number is an estimate of --- and what it is not an estimate of. \hfill
        \scope{[4]}
  \item The median income column is capped at 15. Give one consequence for the
        \emph{model} and one for the \emph{evaluation}, and say which of the
        two you would raise with the client first. \hfill \scope{[5]}
  \item You are given a stratified split on income category and a naive random
        split of the same data. Both report similar RMSE. Does that mean the
        stratification was unnecessary? Answer yes or no, with a reason.
        \hfill \scope{[4]}
  \item The brief says the output feeds a downstream system that expects a
        price. Name the one thing this fact settles about the problem framing,
        and one thing it leaves open. \hfill \scope{[4]}
  \item In the working loop --- specify, generate, read, test, verify ---
        exactly one step is done by the assistant. Name it, and say what goes
        wrong if a student treats a second step as the assistant's. \hfill
        \scope{[3]}
\end{enumerate}

\section*{Summary}
\addcontentsline{toc}{section}{Summary}

Machine learning is worth reaching for when the hand-written rule is long,
unstable or unknown, and it fails silently, which is why evidence rather than
execution is the standard throughout this course. The working method is five
steps of which four are yours: specify with a constraint and a check, read as a
reviewer, test against a case whose answer you already know, and verify that
the number means what it appears to.

The brief was framed before any code: supervised, multiple univariate
regression, batch --- and the assumption that would have made it a
classification problem was checked with the downstream team rather than
guessed. The data was looked at once, in full, which surfaced 207 incomplete
districts, a categorical attribute with a five-district level, a target capped
at \$500{,}001 for 992 districts, scales spanning four orders of magnitude, and
pervasive right skew.

Then the split, \emph{before} the exploration, stratified on the predictor the
experts named --- and measured at 0.36\% worst-case stratum error against
6.4\% for a random split. Only then the exploration, on the training set alone:
income dominates the correlations at 0.688, a default colormap turned out to be
indefensible under a one-line test, three asserted survey artefacts reduced to
one under counting, and a ratio of two weak columns beat both.

\end{document}
